A seguir é apresentada a introdução do trabalho, que situa o tema, justifica sua relevância, delimita os objetivos e descreve a metodologia utilizada.

\section{Contextualização do tema}
Segundo \textcite{carmody_fao2023}, a gestão eficiente dos recursos hídricos é uma questão de importância crescente em um mundo onde a demanda por água aumenta em paralelo com a escassez de recursos naturais. No contexto da agricultura, a irrigação desempenha um papel importante na garantia da produtividade e da sustentabilidade das culturas \parencite{Ramos_irrigacao2022}. No entanto, métodos tradicionais de irrigação muitas vezes carecem de precisão e eficiência, levando ao desperdício de água e outros recursos \parencite{Pereira_irrigation2002}.

Neste contexto, a evapotranspiração (ET) é indicada por \textcite{yang_nature2023} como uma das principais variáveis a serem consideradas na gestão eficiente da água. A ET representa a perda de água do solo por evaporação e pela transpiração das plantas. Entender e monitorar a ET é fundamental para determinar as necessidades hídricas das culturas e otimizar o uso da água na agricultura \parencite{carmody_fao2023, yang_nature2023}.

O desenvolvimento de sistemas de monitoramento em tempo real para a irrigação surge como uma abordagem promissora para otimizar o uso da água e melhorar o rendimento das culturas \parencite{Vijh_system2024, Dutta_system2024, Amine_system2024}. De acordo com \textcite{Prathibha_system2017}, tais sistemas fornecem aos agricultores informações precisas e atualizadas sobre as condições do solo e das plantas, permitindo uma gestão mais eficiente da irrigação ao longo do tempo.

\section{Justificativa}
A justificativa para este estudo reside na necessidade premente de soluções inovadoras que abordem os desafios enfrentados pela agricultura moderna, incluindo a pressão sobre os recursos hídricos e a busca por práticas sustentáveis \parencite{carmody_fao2023}. Ao fornecer aos agricultores ferramentas que lhes permitam monitorar e controlar o processo de irrigação de forma mais precisa e eficiente, espera-se contribuir para a redução do desperdício de água. Isso, por sua vez, pode resultar na otimização do uso de recursos e no aumento da produtividade agrícola.

\section{Objetivos}
Os objetivos deste trabalho incluem o desenvolvimento de um sistema de monitoramento em tempo real que seja capaz de coletar dados relevantes sobre as condições do solo, do clima e das plantas, a fim de fornecer informações para a tomada de decisões relacionadas à irrigação. Além disso, pretende-se avaliar a eficácia e a viabilidade prática do sistema por meio de estudos de caso e análises de desempenho em condições reais de campo.

\section{Metodologia}
A metodologia adotada compreende a pesquisa bibliográfica, o desenvolvimento de protótipos de \textit{hardware} e \textit{software}, bem como testes e validações em campo. Este estudo se limita à implementação de um sistema de monitoramento em tempo real para a irrigação e não aborda questões relacionadas à automação total do processo de irrigação.

\section{Estrutura do trabalho}
Este trabalho está estruturado em cinco capítulos. O Capítulo 1 apresenta esta introdução, contextualizando o tema, justificando sua relevância, delineando os objetivos e descrevendo a metodologia adotada. O Capítulo 2 revisa a literatura relacionada, fornecendo uma base teórica para o desenvolvimento do sistema proposto. O Capítulo 3 descreve em detalhes a metodologia de desenvolvimento do sistema. O Capítulo 4 apresenta os resultados dos testes e validações realizados em campo. Por fim, o Capítulo 5 conclui o trabalho, discutindo suas contribuições, limitações e possíveis direções futuras.

O estudo visa contribuir para o avanço do conhecimento na área da agricultura inteligente, fornecendo uma solução prática e eficaz para o monitoramento em tempo real da irrigação. Essa solução tem potencial para melhorar a sustentabilidade e a eficiência dos sistemas de produção agrícola.